\scanpage{142}
\mspage{87}
% archive sheet number: 141

Une première question, c'est de voir que c'est une équivalence. On voit que c'est ess. surjectif, \unsure{grâces|grâce} à l'existence de $C(X),i_X$. C'est peu! On va essayer de définir une flèche en sens inverse

\[
\bar C:\mathcal M_{/S}W_S^{-1}\unsure{\mathrel{\relbar\joinrel\relbar\joinrel\vert}|\longrightarrow}(\underline{\mathrm{Fib}}\,S)W_S^{-1},
\]

i.e. un \uline{foncteur}

\[
\gamma:\mathcal M_{/S}\longrightarrow(\underline{\mathrm{Fib}}\,S)W_S^{-1}
\]

qui transforme $W_S\subset\mathrm{Fl}\,(\mathcal M_{/S})$ en iso. On choisira, pour $X$ \unsure{ds|dans|de} $\mathcal M_{/S}$, un $C(X)$ dans $\underline{\mathrm{Fib}}\,S$ et un

\[
i_X:X\longrightarrow C(X)\qquad \unsure{v_X|i_X}\in W_S,
\]

et on posera

\[
\gamma(X)=C(X).
\]

La difficulté, c'est de montrer que ça ne dépend pas des choix faits, à isom can. près, et que c'est fonctoriel en $X$ -- et le deuxième implique le premier.

Pour voir que ça marche, visiblement il faut supposer quelque chose de plus. Je vois deux hypothèses qui, l'une et l'autre, \unsure{pourraient} faire marcher: l'une, c'est qu'on \deleted{\unsure{se}} donne déjà $C$ comme foncteur \inserted{et $i_X$ fonctoriel en $X$} (c'est ce qui était le cas dans $\mathrm{Cat}$ \unsure{tantôt|de fait}). L'autre, c'est que
